% ED26-SRC-063-2004: Parsons, complete Fordham manuscript chapter4.
% MS2007 pp.157--280 is distinct from final2008 and the medieval originals.
\chapter{Medieval Supposition Theories: Terms, Context, and Inferential Scope}
\label{sec:p13-medieval-supposition-theories}
\chaptermark{Medieval Supposition Theories}
\index{supposition theory}\index{Parsons, Terence}
\index{ampliation}\index{restriction}\index{descent and ascent}

Supposition theory asks what an occurrence of a term stands for in a
proposition and what inferences its context permits. It is not simply a
list of theories of universals. Terence Parsons's account follows the
development from the later twelfth through the fourteenth centuries,
distinguishing an early quantitative theory from a later theory of
ascent and descent. His complete Fordham manuscript chapter is the
source of this presentation
\cite[pp.~157--280]{Parsons2007FordhamSupposition}.
Its final 2008 publication and the medieval editions and translations it
cites remain distinct sources.

Parsons explicitly offers an ideal framework: differences that matter
for the theory are retained, but no one historical author supplies every
clause of the composite framework. His later repairs using dependent
variables and his four-mode refinement are modern reconstructions,
not recovered medieval symbolic calculi. The starting point is
regimented Latin, whose order gives denoting phrases different scope.
Moving a predicate before the subject can therefore change the assertion
even when the grammatical cases still identify their functions.

\section{Propositions, terms, and the non-vacuous affirmative}
\label{sec:p13-supposition-core}

A \emph{propositio} here is a meaningful sentence, ordinarily a spoken,
written or mental token, rather than the modern abstract proposition
which different sentences might express. Token existence becomes
important in semantic puzzles. Categorematic expressions supply subject
and predicate terms; quantifier signs, the copula and negations are
syncategorematic. A quantifier is not part of the common term it governs.
Singular and common terms are accordingly distinguished from their
complete denoting phrases.

The standard categoricals have singular, universal, particular or
indefinite quantity and affirmative or negative quality. Indefinites
are normally assimilated to particulars in the present core.
Parsons's later-medieval regimentation makes affirmatives with an empty
main term false and negatives with an empty main term true.
This must not be silently replaced by contemporary vacuous universal
truth or by existential import for every particular negative
\cite[pp.~159--165]{Parsons2007FordhamSupposition}.

For a declared extensional comparison let \(S,P\) be current
supposition ranges and write
\[
\begin{array}{ll}
A(S,P):& S\ne\varnothing\ \land\ S\subseteq P,\\
I(S,P):& S\cap P\ne\varnothing,\\
E(S,P):& \neg I(S,P),\\
O(S,P):& \neg A(S,P).
\end{array}
\]
These are modern set clauses for this regimentation, not source-native
medieval glyphs. When \(S\) is empty, A and I are false while E and O
are true. Thus \emph{Some chimera is not a man} can be true because its
subject lacks supposition; it need not assert an existing chimera.
The earlier Abelard treatment distinguishes separative and extinctive
negation and cannot be identified with these clauses without argument.

\begin{hbtheorem}{thm:p13-supposition-square}
 {proved=annex:p13.supposition.square}
 {The square and conversion}
For the four set clauses above, A/O and I/E are contradictory;
A entails I, E entails O, A and E cannot both hold, and I and O
cannot both fail. E and I convert simply; A entails converted I,
and E entails converted O. O does not convert in general, and
does not entail a nonempty subject.
\end{hbtheorem}

The affirmative existence conjunct is essential. The sentence
\emph{Every donkey is an animal} does not become true merely because
there are no donkeys. Conversely, negative truth does not by itself
provide a witness. This is why a modern proof of a source syllogism
must declare its empty-term policy rather than infer one from its name.

\subsection{The syllogistic substrate}
\label{sec:p13-supposition-syllogistic}

Parsons gives twenty-four valid categorical patterns: nineteen named
moods and five unnamed universal-premise patterns with particular
conclusions. His first figure has both direct and indirect forms;
the latter are not silently renamed as his fourth figure. The following
compact transcription preserves the term positions from his table.
Each semicolon separates premises from the final conclusion
\cite[pp.~166--168]{Parsons2007FordhamSupposition}.
\[
\begin{array}{lll}
\text{Barbara}&A(M,P),A(S,M);&A(S,P)\\
\text{unnamed}&A(M,P),A(S,M);&I(S,P)\\
\text{Celarent}&E(M,P),A(S,M);&E(S,P)\\
\text{unnamed}&E(M,P),A(S,M);&O(S,P)\\
\text{Darii}&A(M,P),I(S,M);&I(S,P)\\
\text{Ferio}&E(M,P),I(S,M);&O(S,P)\\[2pt]
\text{Baralipton}&A(M,S),A(P,M);&I(S,P)\\
\text{Celantes}&E(M,S),A(P,M);&E(S,P)\\
\text{unnamed}&E(M,S),A(P,M);&O(S,P)\\
\text{Dabitis}&A(M,S),I(P,M);&I(S,P)\\
\text{Fapesmo}&A(M,S),E(P,M);&O(S,P)\\
\text{Frisesomorum}&I(M,S),E(P,M);&O(S,P)\\[2pt]
\text{Cesare}&E(P,M),A(S,M);&E(S,P)\\
\text{unnamed}&E(P,M),A(S,M);&O(S,P)\\
\text{Camestres}&A(P,M),E(S,M);&E(S,P)\\
\text{unnamed}&A(P,M),E(S,M);&O(S,P)\\
\text{Festino}&E(P,M),I(S,M);&O(S,P)\\
\text{Barocho}&A(P,M),O(S,M);&O(S,P)\\[2pt]
\text{Darapti}&A(M,P),A(M,S);&I(S,P)\\
\text{Felapto}&E(M,P),A(M,S);&O(S,P)\\
\text{Disamis}&I(M,P),A(M,S);&I(S,P)\\
\text{Datisi}&A(M,P),I(M,S);&I(S,P)\\
\text{Bocardo}&O(M,P),A(M,S);&O(S,P)\\
\text{Ferison}&E(M,P),I(M,S);&O(S,P).
\end{array}
\]

\begin{hbtheorem}{thm:p13-supposition-moods}
 {proved=annex:p13.supposition.moods}
 {Every displayed source mood in the declared extension model}
All twenty-four displayed patterns are valid for the stated A/I/E/O
clauses, including their empty-term cases.
\end{hbtheorem}

Mnemonic vowels give the qualities of the premises and conclusion;
initial B/C/D/F indicate reduction to the first four direct moods,
and internal letters indicate simple conversion, conversion by
limitation, premise permutation or indirect proof. These are instructions
for an argument, not extra truth conditions.

For Barocho, \emph{Every man is an animal} and \emph{Some stone is not an
animal} imply \emph{Some stone is not a man}. Assume its contradictory,
\emph{Every stone is a man}. Barbara then yields \emph{Every stone is an
animal}, contradicting the negative premise. Since O is the negation
of A here, the reductio does not require extracting a stone witness
from O alone. The source's conversion comment at printed 168 contains
a repeated particular-to-universal direction; the active deductions
use only the justified conversions specified above.

\section{Syntax and dependent terms}
\label{sec:p13-supposition-syntax}

The extended categorical grammar admits quantifiers on predicates and
negation in more positions. Parsons's minimal grammar begins with the
copula as a partial categorical, allows denoting phrases to be prefixed
until there are two, and allows \emph{not} before a partial categorical.
A complete categorical has two denoting phrases. Moving the copula
past the immediately preceding phrase supplies familiar English order;
Latin case, not surface position alone, identifies subject and object.

Sherwood's equipollences are presented as substitutable phrase patterns:
\[
\begin{array}{ccc}
\text{every A}&\text{no A not}&\text{not some A not}\\
\text{no A}&\text{not some A}&\text{every A not}\\
\text{some A}&\text{not no A}&\text{not every A not}\\
\text{some A not}&\text{not no A not}&\text{not every A}.
\end{array}
\]
Their following context remains in place. They are not identities
between common-term extensions. They change quantifier and negation
placement while preserving the relevant assertion
\cite[pp.~168--173]{Parsons2007FordhamSupposition}.

An explicit modern notation clarifies the distinction:
\[
\begin{array}{ll}
\text{Some donkey not an animal is}:&
 \exists x\in D\ \neg\exists y\in A\ (x=y),\\
\text{Some donkey an animal not is}:&
 \exists x\in D\ \exists y\in A\ (x\ne y).
\end{array}
\]
The first says that some donkey is no animal; the second can merely
compare a donkey with a different animal. These translations are
restricted here to nonempty main ranges. The empty-term regimentation
must be handled separately; it cannot be recovered by inserting an
ordinary existential quantifier into every source particular negative.

The source measures quality by the parity of negating signs, including
\emph{no}, not by an infinitizing prefix such as \emph{non-}. It also treats
conjunction and inclusive disjunction, while distinguishing occasional
exclusive usage; conditionals ordinarily involve necessity of the
connection, and causal, temporal and locational constructions have
different relations. A causal \emph{because} need not denote efficient
causation. None of these is identified merely by notation with a
truth-functional implication.


The account of consequence is modal before it is formal: the premises
cannot be true with a false conclusion. A connection can therefore be
good because of the natures of its terms, as when a donkey's being grey
implies an animal's being grey. Buridan's formal consequence adds
preservation under uniform changes of the terms. Parsons uses modern
validity language for this latter comparison in the nonmodal fragment;
he does not erase the distinction between modal consequence and formal
pattern. The word \emph{consequence} can name either an argument or
its associated conditional. The standard necessary connection should
also be distinguished from an \emph{ut nunc} connection assessed under
current circumstances
\cite[pp.~171--173]{Parsons2007FordhamSupposition}.

The consequence principles reported include transitivity, preservation
of truth and possibility, consequence from an impossibility, consequence
of a necessity from any premise, and deletion of a necessary auxiliary
premise. Detailed consequence theory is outside this source chapter's
scope. So are independent treatments of obligations and insolubles;
syncategorematic doctrines appear only where supposition requires them.
The emerging twelfth-century texts share doctrines and examples whose
common origin the source leaves unknown. Later access to Aristotle's
fallacy discussions helps explain the setting, but does not establish
an identified inventor of the whole theory
\cite[pp.~157--158,172--173]{Parsons2007FordhamSupposition}.

\subsection{From verbs to substantivated terms}
\label{sec:p13-supposition-parasitic}

An adjective can function as a common noun: a \emph{grey} is a grey
thing. An intransitive verb is analyzed into copula and participle;
\emph{Brownie runs} becomes \emph{Brownie is running}, where the participle
is substantivated. This is the regimented analysis, not the aspect
of the English progressive. Existential \emph{A man is} is analyzed
as \emph{A man is a being}, whose predicate ranges over current beings.

For a transitive verb, however, \emph{seeing} depends on the object seen.
Parsons makes that dependency explicit with a \emph{parasitic term}
such as \emph{seeing-\(y\)-thing}. His reconstruction of
\emph{Cicero sees every donkey} has the shape
\[
(\text{Cicero }x)(\text{every donkey }y)
(\text{some seeing-}y\text{-thing }z)\ x=z.
\]
The object phrase binds the variable within the seeing-term. Omitting
that dependency conceals the relation, even if selected examples
continue to look plausible. Parsons does not report a complete medieval
variable calculus for this repair
\cite[pp.~174--177]{Parsons2007FordhamSupposition}.

Intersective modifications combine extensions: a grey donkey is
in the intersection of grey things and donkeys. A relative clause
instead supplies its own predication and tense. \emph{Every white man
runs} restricts the men under discussion; a nonrestrictive assertion
that every man, who is white, runs adds a different commitment.
Deleting \emph{which is} from an appropriate relative clause preserves
its concurrent-tense reading, but deleting an explicitly past relative
verb does not in general preserve its distinct time.

A genitive also creates a dependency. In
\emph{Of every woman a donkey is lame} each woman can have her own
donkey; \emph{A donkey of every woman is lame} may require one donkey
possessed by all of them. A modern expansion uses
\emph{donkey-of-\(x\)}, with possession indexed by the earlier possessor
phrase. Relational \emph{father of} is not the same as the contextual
possession of a dog. Combining \emph{the dog is yours} and \emph{the dog is
a father} does not make it your father
\cite[pp.~179--184]{Parsons2007FordhamSupposition}.


The relative clause can be represented by a proposition with a gap:
an object belongs to the complex term precisely when it belongs to
the head term and satisfies the clause with that object in the gap.
The relative pronoun takes the grammatical case appropriate to the
missing phrase. A bare possessive such as \emph{Brownie is Socrates's}
can be completed with an unspoken \emph{thing}. This also explains a
following possessor: \emph{Some donkey of every farmer is brown}
can be expanded to a donkey which is every farmer's thing, then reduced
by deleting the relative pronoun and copula. A preceding possessor
instead forms an independent denoting phrase that binds the possessed
term. These are Parsons's reconstruction devices, not a claim that
medieval texts explicitly used gaps and bound variables
\cite[pp.~180--185]{Parsons2007FordhamSupposition}.

\begin{hbtheorem}{thm:p13-supposition-scope}
 {proved=annex:p13.supposition.scope}
 {Scope and binding}
In the declared nonempty relational comparison, each person seeing
some donkey need not yield one donkey seen by every person.
Likewise a girl giving every present to some boy need not give some
one boy every present. Possessing an object that is a father need
not make that object the possessor's father.
\end{hbtheorem}

\subsection{Infinitizing negation and exponibles}
\label{sec:p13-supposition-infinitizing}

\emph{Non-man} changes the common term; it is not the same operation
as negating the proposition. Its extension is sometimes confined to
beings and sometimes taken more widely. The current present-tense
core does not settle all such ontological alternatives. In particular,
contraposition can fail because the complemented subject becomes
empty: \emph{Every man is a being} can be true while
\emph{Every non-being is a non-man} is false.

Obversion from an affirmative to an appropriate negative is sound
in the declared extension core, but its converse can introduce
an unwarranted affirmative existence commitment. A negative chimera
sentence does not establish an affirmative non-man chimera
\cite[pp.~177--179]{Parsons2007FordhamSupposition}.

\begin{hbtheorem}{thm:p13-supposition-complements}
 {proved=annex:p13.supposition.complements}
 {The four obversions and the empty-complement obstruction}
Taking complements relative to a fixed domain, the four implications
\(A(S,P)\Rightarrow E(S,\overline P)\),
\(A(S,\overline P)\Rightarrow E(S,P)\),
\(I(S,P)\Rightarrow O(S,\overline P)\), and
\(I(S,\overline P)\Rightarrow O(S,P)\) hold.
Their unrestricted converses and unrestricted contraposition fail.
\end{hbtheorem}

An \emph{exponible} is explained by several simpler assertions.
\emph{Only a man runs} requires a running man and excludes non-men
from running; \emph{No man except Socrates runs} combines Socrates's
running with exclusion of every other man. Beginning and ceasing
combine present and past affirmative and negative claims.
The analysis supplies logical structure rather than one opaque term,
but a term can acquire different modes in different exponents.
That issue returns below; it prevents a single undifferentiated
mode from being assigned mechanically
\cite[pp.~184--186,259--260]{Parsons2007FordhamSupposition}.


\section{Signification, supposition, and kinds of reference}
\label{sec:p13-supposition-signification}

Signification belongs to an expression before its occurrence in a
particular proposition; supposition belongs to that occurrence. Parsons
sets out two backgrounds for this distinction. In a realist account a
common word is conventionally imposed on a form, and the form is shared
by individuals. The word thereby also signifies those individuals. In a
nominalist account the intermediary is a concept naturally of individuals,
rather than a universal outside the mind. The two diagrams have similar
logical organization without having the same ontology. Ockham's wording
is more precise than a simple chain of indirect reference: spoken words
directly, though secondarily, signify what the associated concepts signify
\cite[pp.~186--188]{Parsons2007FordhamSupposition}.

The familiar three kinds of supposition distinguish the uses of
\emph{donkey} in \emph{Donkey is bisyllabic}, \emph{Donkey is a
species}, and \emph{A donkey is running}. Material supposition concerns
an expression or a related expression; simple supposition concerns the
associated form or concept; personal supposition concerns the ordinary
significata. ``Personal'' therefore includes stones and other nonpersons.
Nor can the kinds be defined just by the category of the object referred
to. In \emph{Some words are nouns}, \emph{word} personally supposits
for words, including itself; in \emph{Every species has a genus},
\emph{species} personally supposits for species. The way the occurrence
signifies, rather than self-reference or reference to a form alone,
matters to its classification.

The authors do not apply this taxonomy uniformly. Peter of Spain gives
predicates simple supposition, even the predicate of an ordinary
affirmative categorical. Sherwood instead speaks of their appellation.
Later nominalists generally allow subject and predicate alike to
supposit personally. The extension clauses of
\cref{sec:p13-supposition-core} represent that later treatment. Their
historical purpose is in part to account for truth without extra-mental
universals. They are not a complete recursively specified semantics:
universally quantified predicates require further clauses, and
\emph{Some A sees every B} introduces a relation or a third, dependent
term after analysis into copula and participle
\cite[pp.~188--189,192--194]{Parsons2007FordhamSupposition}.

\subsection{Mental language and connotation}
\label{sec:p13-supposition-mental}

For Ockham, concepts are also words of mental language. Their
signification is natural; the relation of spoken or written expressions
to them is conventional. Mental propositions supply the primary truth
bearers to which conventional sentences are subordinated. The proposal
requires those grammatical distinctions that can change truth: case
distinguishes who sees whom, and number distinguishes being an animal
from being the animals. A declensional difference that contributes no
such distinction need not occur in mental language. The same economy
motivates the rejection of synonymous mental expressions. These are
claims of the historical theory, not an empirical demonstration that all
human thought uses one language
\cite[pp.~189--190]{Parsons2007FordhamSupposition}.

Ockham's absolute/connotative distinction raises a separate issue. A
connotative term such as \emph{white} primarily signifies white things
and secondarily signifies whitenesses; its nominal definition uses a
construction such as a thing having whiteness. An absolute term has a
real or quidditative definition, a characterization of the things, which
need not be synonymous with the term. Several true characterizations of
human beings need not give the same account of the word \emph{man}.
Buridan associates a strict nominal definition with a complex concept.
It does not follow that Ockham eliminates every connotative term into
absolute vocabulary. Parsons leaves both that reduction and the exact
relation between a connotative term and its nominal definition disputed.
If a simple connotative mental term and its complex definition both
exist, the prohibition of mental synonyms prevents their being treated
as two synonymous mental expressions
\cite[pp.~190--192]{Parsons2007FordhamSupposition}.

\subsection{Material supposition and the limits of quotation}
\label{sec:p13-supposition-material}

The background to the three kinds includes the older discussion of
\emph{univocation}. Boethius's apparently opposed claims that man walks
and man does not walk can take the term for an individual in one case
and the species in the other while retaining its definition. This was
distinguished from equivocation, in which a common spelling has different
definitions. Other writers disputed the boundary; Buridan can classify
the change as equivocation because a different mental concept governs
the use. No single modern meaning of ``univocal'' should be read back
into all these accounts
\cite[pp.~192--193]{Parsons2007FordhamSupposition}.

Material supposition supplies resources that quotation often supplies
today, but the two are not identical. It can connect a written token
with a spoken token, tokens of a type, an inflected relative, or an
expression taken under another supposition. Burley's taxonomy includes
spoken-for-spoken and written-for-spoken reference, a clause standing
for its corresponding indicative sentence rather than itself, and a
word standing for its own personal use. A word which cannot function
substantivally in one setting can itself be mentioned substantivally in
another. Later use of the imported article \emph{ly} makes such mention
explicit; this does not make quotation the original analysis of every
case \cite[pp.~194--196]{Parsons2007FordhamSupposition}.

Parsons's naming example illustrates a forced material context.
\emph{I am called Herman} and \emph{Herman is my name} do not license
\emph{I am called my name} on the same intended reading. In the second
premise \emph{my name} denotes Herman; after \emph{am called} the
phrase can instead be taken for its own words. The apparent identity
substitution changes the way the inserted expression is used. Adding
quotation marks mechanically does not repair a construction that itself
forces mention. Material reporting can also leave open whether the
reported speaker used some words personally or materially. Explicitly
nested quotation marks specify that difference; a report of the words
alone need not specify it
\cite[pp.~196--197]{Parsons2007FordhamSupposition}.

\subsection{Three cases of simple supposition}
\label{sec:p13-supposition-simple}

In the first and most straightforward case, \emph{Man is a species}
takes the term for its form, or for its associated concept on a
nominalist interpretation. An early alternative treats this occurrence
as a discrete name of the species. Ockham retains simple supposition for
a mental intention, used nonsignificatively. Buridan instead calls this
material supposition: a concept is a mental word, so no third kind is
needed alongside material and personal. That is a substantive
reclassification, not the discovery that all earlier authors used two
kinds \cite[pp.~197--199]{Parsons2007FordhamSupposition}.

Sherwood's second case connects the form with individuals in which it
is preserved. His example of man as the noblest creature permits a
claim about Socrates \emph{as a man}, without making Socrates simply
the noblest creature. The reduplicative qualification does essential
work. Parsons supplies no complete semantics for it, and the present
presentation does not invent one. Ockham rejects the proposed simple
supposition: the literal proposition is false on a personal reading,
but its intended assertion can be that every human is nobler than every
nonhuman corporeal creature. The restriction to corporeal creatures
matters. Other comparative and superlative examples show why the
unworked reduplicative analysis cannot be generalized without further
argument \cite[pp.~199--201]{Parsons2007FordhamSupposition}.

The third, unfixed use concerns propositions such as pepper's being
sold here and in Rome, the displayed plant's growing here and in the
garden, and woman having both damned and saved humankind. The intended
individuals differ across the predications. They are neither one
unchanging individual nor a species that is literally sold or grows.
Sherwood relates the form to undesignated individuals. Ockham and
Burley prefer an intended paraphrase: a plant \emph{of this sort} grows
in the garden. A literal demonstrative still denotes the displayed
plant and may make the original assertion false. These examples expose
the difference between charitable paraphrase and a new semantic kind
\cite[pp.~201--202]{Parsons2007FordhamSupposition}.

\section{Ampliation, restriction, and time}
\label{sec:p13-supposition-ampliation}

In the default nonmodal present tense, a term personally supposits for
its presently existing significata. Thus a present affirmative about
the deceased Socrates fails through an empty subject on this account.
Tense, modality, mental vocabulary, and semantic vocabulary can enlarge
the range. This is ampliation. A modifier can instead restrict a range.
The terminology takes the present range as its reference point; it does
not establish that every ontology contains past, future, merely possible,
or impossible things
\cite[pp.~202--205]{Parsons2007FordhamSupposition}.

In \emph{Some bishop was running}, the subject can concern a present
bishop who formerly ran or a former bishop who formerly ran. Ockham
distinguishes readings; Buridan gives disjunctive truth conditions.
Parsons follows the latter for exposition. The predicate concerns past
runners, not merely present runners. The future case is analogous.
The copula in the resulting comparison expresses tenseless identity;
the tense changes the ranges compared. This is why a past sentence is
not generally equivalent to saying that its present-tense wording was
once true. A now-white object may formerly have been black without
ever having been both white and black
\cite[pp.~206--208]{Parsons2007FordhamSupposition}.

\begin{hbtheorem}{thm:p13-supposition-time}
 {proved=annex:p13.supposition.time}
 {Independent temporal ranges do not synchronize their witnesses}
In a finite comparison with a present time and two past times, the
disjunctive subject-range account can make \emph{A white was black}
true although \emph{A white is black} was never true. It can also make
\emph{A bishop was grey} true without any time at which one individual
was simultaneously a bishop and grey.
\end{hbtheorem}

This last possibility is a theoretical boundary of the articulated
account. It separately allows a past time for being a bishop and a past
time for being grey. A reading with one shared past time is a stronger
condition. Examples such as a virgin's having been pregnant or an old
man's being a boy in the future were used to display this lack of
coordination; their intelligibility depends on the temporal readings,
not on one individual's being both at once. Parsons does not settle
whether every ordinary-language use demands the stronger reading.

Relative clauses contribute their own tense. In \emph{A donkey which
will be grey was brown}, the proposed analysis ampliates the head noun
to present or past donkeys, retains the relative clause's future grey
condition, and uses past brown things for the main predicate. This
does not force the times within a past relative clause and the main
clause to coincide. Possessive dependent terms raise another problem:
the time of being a donkey and the time of belonging to its owner both
matter, but the source identifies no worked-out coordination rule.
Buridan's requirement that eating an apple in Paris happen at a time
when it is an apple and the eater is in Paris is natural; independent
ampliation of each term does not yet implement it
\cite[pp.~208--210]{Parsons2007FordhamSupposition}.

Embedded speech offers a different resource. A past complaint that
Plato had hit Socrates can be understood through a mental proposition
existing at the time of the complaint, whose own past tense points
earlier than that time. This yields a past-under-past reading. A report
that Socrates thought Euthyphro was a fool can instead report a
present-tense content at the time of thinking. The source suggests
treating the subordinate past as present in that case, but does not
give a full sequence-of-tense theory
\cite[pp.~210--211]{Parsons2007FordhamSupposition}.

Restriction by modification should be separated from these temporal
issues. All parties agree that \emph{grey donkey} supposits for grey
donkeys. Some additionally say that \emph{grey} restricts the simple
term \emph{donkey} itself, and vice versa. For the truth of a sentence
whose main term is the whole complex, both analyses take the same
intersection. A modifier such as risible can leave the complex term
empty. Calling that alienation does not by itself alter the sentence's
declared truth condition
\cite[pp.~205--206]{Parsons2007FordhamSupposition}.

\subsection{Modal, mental, and semantic ampliation}
\label{sec:p13-supposition-modal}

Buridan distinguishes three constructions. An assertoric sentence can
contain a modal noun or adjective, as when a proposition is called
possible. A composite modal attributes modality to a dictum, what is
said by a clause. A divided or genuinely modal construction modifies
the copula: an animal can be grey, or necessarily is an animal. Modal
adverbs, the auxiliary \emph{can}, and Latin modal-adjective constructions
with accusative subjects and infinitives supply different grammatical
forms. Parsons's own view that the last may be a raised-subject
construction is explicitly not Buridan's classification
\cite[pp.~211--213]{Parsons2007FordhamSupposition}.

The semantic option Parsons develops treats these latter constructions
on the temporal pattern. The subject range includes possible things;
the predicate range comprises possible or necessary instances according
to the mode. The resulting identity comparison does not by itself tell
us which things can be donkeys, which are necessarily donkeys, or whether
never-actual possible things exist. Those depend on metaphysics. A
composite necessity claim about a universal affirmative also carries
the affirmative's existence condition: the possible absence of donkeys
can obstruct necessity even when being a donkey entails being an animal.
No unrestricted substitution of a modern possible-world semantics for
these range clauses is made here
\cite[pp.~213--214]{Parsons2007FordhamSupposition}.

Words pertaining to the soul can allow their affected terms to range
over everything they signify, including impossible objects if the
ontology admits them. An opinable Antichrist can be future; thinking
of a golden mountain need not require a present mountain. A mental verb
need not ampliate its preceding subject: a present assertion that the
deceased Socrates thinks is still false on the default restriction.
Buridan's rejection of impossible beings makes an affirmative about a
chimera empty even in an ampliating context. An expanded license to
supposit is not a guarantee of an inhabitant
\cite[pp.~203--205,214]{Parsons2007FordhamSupposition}.

Semantic words such as \emph{signify} and \emph{supposit} similarly
permit discussion of all significata, not just present instances.
Assertions made with them still have commitments. The affirmative that
every chimera is signified would commit its speaker to a nonempty range;
Buridan denies the corresponding signification claim rather than
accepting impossible chimeras. The concept can remain meaningful by
composition from meaningful components while having no significata.
The semantics therefore provides ways of stating ontological commitments
without resolving all of them in advance
\cite[p.~215]{Parsons2007FordhamSupposition}.

\section{Nonextensional constructions}
\label{sec:p13-supposition-intensional}

Latin belief reports commonly use an accusative subject with an
infinitive, rather than an English that-clause. The infinitive's time
normally depends on the main verb; a that-clause has a tense of its
own, even when interpreted dependently. Translating away the difference
can obscure the form a logician is analyzing. Parsons compares three
accounts of what such a clause supposits for
\cite[pp.~215--218]{Parsons2007FordhamSupposition}.

For Ockham it can supposit materially for the corresponding indicative
proposition. If that is a mental proposition, the object of belief is
a mental sentence; if spoken or written propositions are meant, the
objects differ accordingly. Abelard's non-things and Wodeham's objects
signifiable by a complex provide a different option. ``Not a thing''
here denies membership in the Aristotelian categories; it is not simply
a denial that anything is signified. Wodeham distinguishes taking the
clause for a linguistic dictum from taking it for what the dictum
signifies. These views must not be collapsed into Ockham's material
supposition or treated as an established modern abstract-proposition
ontology.

Buridan offers an account on which the clause supposits for its subject's
suppositum when the predication is correct, and otherwise for nothing.
Socrates-to-love-God is Socrates if he loves God, but is empty otherwise;
a past hatred can concern the same person under past ampliation. Parsons
also gives the unaddressed challenge: if God loves Socrates while
Socrates hates God, identity appears to make God love Socrates-to-hate-God.
The source does not supply a solution, so this is retained as a limit
of the account rather than advertised as an extensional semantics of
all belief reports \cite[pp.~217--218]{Parsons2007FordhamSupposition}.

An intensional transitive raises a related difficulty. One can know
one's father without recognizing the approaching person as one's father.
Buridan and followers appeal to a \emph{ratio}: knowledge of an object
under one presentation does not imply knowledge under another. Identity
of the object preserves \emph{knowing the approaching person under the
ratio of one's father}; it does not preserve the omitted ratio in
\emph{knowing the approaching person}. Fronting the object can instead
express knowing it under some ratio. A concept is a natural candidate
for the ratio, but Parsons does not report a complete theory of it
\cite[pp.~218--219]{Parsons2007FordhamSupposition}.

The promise or debt of a horse tests another failure of unrestricted
descent. A promise to supply a horse need not promise this horse, nor
any separately specified horse. Realists can say that the universal is
owed and is delivered through an instance; Wycliff makes \emph{promise}
induce simple supposition in the following common term. A singular term
is personal instead, leaving the inference from a particular promise
to the general promise in need of explanation. Buridan argues by parity
for a promise of every horse rather than none. Ockham takes the promise
to concern bringing about the truth of giving a horse, without requiring
any one singular giving proposition. Buridan's objection about every
or no way of discharging it does not answer the possibility that any
one of several ways suffices. The source records the disagreement; it
does not infer a determinate owed horse from the grammatical object
\cite[pp.~220--221]{Parsons2007FordhamSupposition}.

\section{The early quantitative theory of modes}
\label{sec:p13-supposition-early-modes}

Discrete personal supposition belongs to a whole singular term, such
as a proper name, a demonstrative, or \emph{this donkey}. The modes of
common personal supposition classify occurrences of common terms as
affected by their governing signs. They do not classify the word-type
once for all, nor simply replace its range with one selected object.
Some writers extend modes to common material uses, as in counting
occurrences of a word; that extension is not assumed throughout the
tradition \cite[pp.~222--225]{Parsons2007FordhamSupposition}.

For the standard categoricals the common classification is:
\[
\begin{array}{c|cc}
 &\text{subject}&\text{predicate}\\\hline
A&\text{distributive}&\text{merely confused}\\
I&\text{determinate}&\text{determinate}\\
E&\text{distributive}&\text{distributive}\\
O&\text{determinate}&\text{distributive}
\end{array}
\]
Peter of Spain's predicates are an important exception because he
classifies them as simple. The usual division first separates
determinate from confused, then divides confused into distributive and
merely confused. Confusion in this terminology concerns taking things
together; it does not mean a mistake by the reasoner.

Sherwood, Lambert, and Peter give quantitative accounts of modes. Their
formulations differ. Parsons reconstructs determinate supposition
by a proposition's being true for at least one suppositum. Replacing the
whole governing denoting phrase by a variable for that suppositum
explains ``true for.'' The interpretation is existential, not uniquely
existential: several people may run while \emph{A man runs} remains
determinate. Peter explicitly distinguishes what the term stands for
from what makes this occurrence's proposition true. Lambert's wording
mentions sufficiency; taking it as sufficiency alone would wrongly make
the predicate of a universal affirmative determinate. Parsons therefore
reads the intended condition as an equivalence
\cite[pp.~225--227]{Parsons2007FordhamSupposition}.

For distributivity, ``many'' cannot require several distinct witnesses:
every person might see the same person. On Parsons's reconstruction it
excludes the determinate pattern, and ``any'' says that the original
proposition entails its instance for each suppositum. This is a
one-way condition. The predicate of a particular negative requires that
asymmetry: some dog being outside the spotted things entails some dog
being distinct from each specified spotted thing, but the collection
of these claims need not supply a dog outside the entire class. Merely
confused supposition is the remaining personal case, neither
determinate nor distributive. The reconstruction makes the predicate
of a universal affirmative its paradigm
\cite[pp.~227--230]{Parsons2007FordhamSupposition}.

These semantic readings and the proposed causes of modes are different
parts of the theory. Peter requires a universal sign for distribution
and denies that negation itself is such a sign. Sherwood's rule for a
distributive sign gives distribution to the adjacent term; an
affirmative sign confuses the remote term, while a negative sign
distributes it. Parsons supplements these with a mobility rule and an
immobility default: distribution permits descent unless impeded, and
a term is immobile unless something makes it mobile. In applying those
supplements, the rule permitting mobility is the former, not the latter;
the manuscript's occasional reference to the immobility rule cannot
justify the inference
\cite[pp.~229,233--236]{Parsons2007FordhamSupposition}.

\subsection{Locality, mobility, and failed inference rules}
\label{sec:p13-supposition-early-rules}

Mobility permits replacement by any singular inferior, with a negative
governing sign's negation retained. Thus a donkey's being distributed
in \emph{Every donkey runs} licenses descent to Brownie, given that
Brownie is a donkey. Sherwood and Lambert let the mode remain locally
distributive when \emph{only} blocks that descent. They likewise retain
local distribution in \emph{Not every man runs}. Peter can allow
mobility without distribution: the inference from \emph{Not some man
is just} to Socrates's not being just follows by contraposing the
ordinary particular ascent, even though negation is not a distributing
sign for him. Hence mobility and distribution cannot be interchanged
in these early accounts
\cite[pp.~231--237]{Parsons2007FordhamSupposition}.

The local account also explains why an equipollence can change mode
while preserving consequence. \emph{No donkey is a stone} and
\emph{Not some donkey is a stone} need not have the same local
distribution. But even within extended categoricals the rules are
incomplete. The predicate in \emph{Every A is every B} receives two
influences whose compatibility depends on whether ``merely confuses''
means only adding confusion or imposing exclusive mere confusion.
In \emph{No A is every B} the global inferential test makes B merely
confused, although both signs seem to distribute it by the early rules.
In \emph{Some A is not every B} and \emph{Some A is not no B}, the
test makes B determinate, again conflicting with the unrevised causes.
The table of the early theory is therefore not a general scope
algorithm \cite[pp.~230--231,236--237]{Parsons2007FordhamSupposition}.

Sherwood's remaining rules warn against moving from mere confusion
to distribution, from many separate determinates to one shared
determinate, or from a determinate to distribution. They allow a move
from determinate to merely confused and from distributive to determinate
in their illustrated circumstances. Self-seeing people provide the
standard countermodel: everyone sees someone, but no one person is
seen by everyone. Separate witnesses for Socrates, Plato, and the rest
can yield an existential under a universal, not one existential outside
it. Likewise, someone's not being seen by Socrates does not imply that
Socrates sees nobody
\cite[pp.~238--239]{Parsons2007FordhamSupposition}.

Three restrictions prevent these insights from becoming unrestricted
rules. First, embedding a valid implication in the antecedent of a
conditional reverses its permissible direction; unchanged local modes
do not detect that. Second, altering one sign can also alter another
term's mode: replacing \emph{Every donkey is an animal} by
\emph{Some donkey is every animal} changes both roles. Third, even
unchanged traditional modes do not determine dependence: giving every
present to some boy need not give one boy every present. The finite
countermodels in \cref{thm:p13-supposition-scope} establish the last
failure. Parsons's restricted proposals consequently hold the remaining
wording, grammatical relations, and relevant scopes fixed; they are
not a license for arbitrary rearrangement
\cite[pp.~239--241]{Parsons2007FordhamSupposition}.

\section{The later theory of ascent and descent}
\label{sec:p13-supposition-later-modes}

The later account changes both the criterion and its scope. Burley,
Ockham, Buridan, and their successors use admissible inferences to
classify terms in a whole proposition. Embedding can change the mode,
as well as impede or enable descent. The principal composite account
in Parsons is deliberately assembled from their best-fitting clauses;
no one writer states precisely the entire composite
\cite[pp.~241--245,260--266]{Parsons2007FordhamSupposition}.

To state the tests, fix a term occurrence in a proposition and a list
exhausting its supposita. An instance replaces its governing denoting
phrase by a singular demonstrative for a listed suppositum. Negation
carried by a negative governing sign must be retained. \emph{Descent}
passes from the original proposition to instances; \emph{ascent} passes
back. Exhaustiveness and correct reference are accompanying information,
not facts inferred merely from writing a list. Under the composite
definitions:

\begin{enumerate}
\item A determinate occurrence permits descent to the disjunction of
  its instances and ascent from any one instance.
\item A distributive occurrence permits descent to their conjunction,
  but not ascent from an arbitrary single instance.
\item A merely confused occurrence permits neither conjunctive nor
  disjunctive propositional descent, but permits ascent from an instance.
\end{enumerate}

For \emph{Some donkey is spotted}, an exhaustive disjunction says that
this donkey is spotted or that donkey is spotted, and so on. Each
disjunct also yields the original. For \emph{Every donkey is spotted},
the corresponding conjunction follows, but one spotted donkey does
not suffice for ascent. For the predicate of \emph{Every donkey is
a mammal}, neither a conjunction nor a disjunction of claims that every
donkey is this one mammal follows. Any such uniform identity claim
does imply the original. The predicate is therefore merely confused.

The negative forms preserve the contrast. \emph{Some stone is not
a donkey} implies, for each specified donkey, that some stone is not
that donkey. The witnesses in the separate conclusions can differ.
Thus this descent does not assert that the original O-proposition can
be recovered from their conjunction. In \emph{Not every man runs},
the later global test makes the subject determinate, unlike Sherwood's
local distribution: with a nonempty range it is a disjunction of
nonrunning instances. Empty ranges require the qualification below
\cite[pp.~242--246]{Parsons2007FordhamSupposition}.

\subsection{What the tests presuppose}
\label{sec:p13-supposition-descent-limits}

An occurrence, not a spelling, is tested. \emph{Every donkey is a
donkey} has a distributed subject and a merely confused predicate in
the intended formal classification. Applying arbitrary consequence to
that very sentence would let the repeated noun create accidental
entailments. Parsons therefore replaces its other occurrence by an
independent schematic term when testing one occurrence. He proposes
extending this discipline inside complex terms: the contradictory
\emph{donkey which is not a donkey} otherwise makes every consequence
trivial. An inference from a contradiction is not automatically an
instance of the particular ascent rule being studied
\cite[pp.~246--248]{Parsons2007FordhamSupposition}.

Empty and singleton ranges also collapse distinctions in the literal
tests. The texts do not provide a settled theory of empty conjunctions
and disjunctions. Ockham's proviso that a relevant term have many
inferiors suggests testing a schematic occurrence under a nonempty,
plural range. This is not a revision of the empty-term truth clauses
in \cref{sec:p13-supposition-core}. It is a restriction on the
diagnostic classification: the truth of an empty-subject O does not
provide the existential witness used by a determinate descent. The
finite checks below explicitly assume nonempty ranges, and use ranges
of at least two things to witness failures of inference
\cite[pp.~248--249]{Parsons2007FordhamSupposition}.

Mobility becomes built into distribution for Buridan; Marsilius likewise
rejects an immobile distributed term. An exceptive restricts the range
to which descent applies, instead of making that restricted range
immobile. Ockham allows descent ``in some way'' after rewording, but
does not specify a general class of permitted rewordings. Burley speaks
of distribution with respect to an exception, although an instance such
as Plato-besides-Socrates is not obtained in the expected way. Paul of
Venice gives modal examples in which descent fails. These attempts to
retain immobile distribution are recorded as variants; they do not
supply one completed alternative test. Parsons's main later account
uses Buridan's mobile distribution
\cite[pp.~249--250]{Parsons2007FordhamSupposition}.

\subsection{Causes of modes and the double-distribution problem}
\label{sec:p13-supposition-later-causes}

The later causal rules are recursive. Determinacy is the default. A
universal affirmative distributes its own term and confuses a following
main term within scope if that term is not already confused. A universal
negative distributes its own term and changes the modes of following
terms in its scope. Negation likewise changes terms in its scope;
it does not affect an earlier term merely because that term occurs in
the same sentence. These rules give the usual A/I/E/O classification
\cite[pp.~250--253]{Parsons2007FordhamSupposition}.

The difficulty lies in a term already distributed by another sign.
Buridan normally makes it merely confused; Ockham makes it determinate.
The two answers are appropriate in different cases. The term governed
by \emph{every} in \emph{Some A is not every B} becomes determinate
under the global test. But in \emph{Not: some farmer sees every
donkey}, the donkey term is merely confused: no one donkey need defeat
every candidate farmer, although a donkey seen by no farmer would
suffice. Buridan separately exempts two negations relating in the same
way to their scope, where they cancel. Neither author's single rule
for all doubly distributed terms captures both sorts of example.

Infinitizing negation presents a further context change. Buridan lets
\emph{man} be distributed in \emph{A non-man runs}, with that
distribution removed in \emph{Every non-man runs}. A descent retains
the prefix and its external sign, producing a non-Socrates term rather
than replacing the whole denoting phrase by Socrates. This differs
from the main-term substitution test above. Extending the test can be
useful, but it is an extension to a non-main occurrence. Peter of Spain
rejects the relevant distributing role of negation; Buridan's attribution
does not remove their disagreement
\cite[pp.~253--254]{Parsons2007FordhamSupposition}.

\subsection{Dependent ranges and restricted descent}
\label{sec:p13-supposition-dependent}

Consider \emph{Some farmer's every donkey runs}. In a modern
regimentation it has an existential farmer followed by a universal
over that farmer's donkeys. The range of the second term is a function
of the first witness. The causal rule predicts distribution, but an
unqualified descent to every donkey in the domain is false. Descent
to every donkey belonging to some farmer is also false. The intended
range is the donkeys of the farmer witnessing the original assertion.
Repeating an independent \emph{some farmer} in each conjunct loses
that shared reference
\cite[pp.~255--256]{Parsons2007FordhamSupposition}.

\begin{hbtheorem}{thm:p13-supposition-dependent-range}
 {proved=annex:p13.supposition.dependent-range}
 {A dependent range cannot be replaced by the unrestricted range}
There is a finite nonempty farm model in which some farmer owns
donkeys and every donkey of that farmer runs, while a donkey owned
by another farmer does not run. There is also a model in which every
animal of a farmer runs, but replacing that farmer's animals by that
farmer's donkeys makes the non-vacuous affirmative false, even though
every donkey is an animal and there are donkeys in the domain.
\end{hbtheorem}

Buridan's proposed minor premise, that Brownie is \emph{that man's}
donkey, points to the right dependency, but the ordinary list test does
not bind the anaphoric phrase. With a fixed singular owner, such as
Socrates, there is a natural extension: list exactly Socrates's donkeys.
With a quantified owner, a dependent witness and a conditional descent
would need an explicit rule. Parsons therefore does not credit parasitic
terms with ordinary modes merely because the causal algorithm assigns
them labels. A main term can have global quantificational import without
passing the given independent ascent/descent tests.

\subsection{Inferiors, superiors, and exponibles}
\label{sec:p13-supposition-monotonicity}

For an independently ranging personal term, the inferential applications
are clear. A superior-to-inferior move in a distributed position uses
an additional premise that every inferior is a superior. In an affirmative
setting that premise supplies the inferior's nonemptiness as well as
the inclusion. Barbara and Celarent are basic examples; seeing every
horse and Brownie's being a horse also entail seeing Brownie. In a
determinate or merely confused position the direction is reversed:
an inferior can be replaced by a superior, again with the relevant
inclusion premise. Darii is an instance. Ockham explicitly restricts
the latter principle to personal, significative uses; material and
other uses furnish counterexamples
\cite[pp.~256--258]{Parsons2007FordhamSupposition}.

Burley's formulation also permits loss of distribution in suitable
conclusions: from every A's being B and every B's being C, infer either
every A's being C or some A's being C. A negative case allows
\emph{Not some D is E}, with every F being E, to yield
\emph{Not some D is every F}. The nonempty F premise prevents the
universal predicate from becoming vacuously true. These examples have
the stated personal ranges and scope. The dependent farm model shows
why the principle cannot simply replace one common head inside a
varying possessive range
\cite[pp.~258--259]{Parsons2007FordhamSupposition}.

Exponibles complicate the assignment of one mode to one visible word.
Beginning to be white expands into a past negative and a present
affirmative; the repeated white term receives different modes in those
exponents. Buridan allows a term to have the supposition or suppositions
found there. Ockham can say that the unexpanded occurrence has none of
the normal modes, but he does not apply that verdict uniformly to
exclusive propositions. In \emph{Only an animal is a man}, the
exponents' subjects differ, and he retains mere confusion in the
original. Parsons's suggested policy is to test the original first and
use divergent exponents to explain a failure. He marks the policy as an
interpretation, not an explicit completed rule of Ockham's
\cite[pp.~259--260]{Parsons2007FordhamSupposition}.

\subsection{Variants of the tests}
\label{sec:p13-supposition-variants}

Burley's determinate definition omits ascent. Taken literally, that
would also capture universal subjects whenever a nonempty conjunction
of instances implies their disjunction. Parsons regards the omission
as an oversight in light of the examples. Ockham and Buridan include
ascent, and Albert's formulation supplies it through the sufficiency
of the disjunction. In the distributive case Buridan omits the
prohibition on single-instance ascent. Whether that prohibition is
redundant in the restricted language requires the discipline of excluding
accidental repeated-term consequences; it is not a theorem here about
arbitrary propositions
\cite[pp.~260--262]{Parsons2007FordhamSupposition}.

Paul's \emph{Logica Magna} instead requires return from the entire
conjunction for mobile distribution. This works for universal subjects
and universal-negative predicates but fails for a particular-negative
predicate. Calling the remaining case immobile is unusual because
individual descent remains possible. The distinction nevertheless
anticipates one part of Parsons's later refinement.

Ockham adds a different test for mere confusion: descent to a
\emph{disjunctive term}, as in every horse being this animal or that
animal. This is not a disjunction of propositions asserting one animal
of every horse. The test depends on the grammatical denoting phrase:
\emph{every} naturally expands conjunctively, while \emph{some}
expands disjunctively. Changing signs by an equipollence can therefore
alter eligibility for this test without altering the relevant
inferential behavior. In particular, \emph{No animal is every man}
cannot in general be replaced by \emph{No animal is this man or that
man}. The former can be true when there are several men, though the
latter fails because men are animals
\cite[pp.~262--265]{Parsons2007FordhamSupposition}.

Dorp proposes moving the verb to the end and negations to its right
before applying the test. That normalizes the remaining main signs to
\emph{every} and \emph{some}, potentially making the extra disjunctive
term condition redundant. Other proposals add a fourth mode for
conjunctive-term descent or allow either term expansion. These are not
Parsons's later split of distribution. Buridan takes another route by
omitting ascent from his merely confused definition; then every
non-determinate, non-distributive occurrence falls into it, including
occurrences inside belief contexts or relative clauses which fail the
composite account's ascent requirement. This creates a broader residual
category rather than the same three-way classification
\cite[pp.~265--266]{Parsons2007FordhamSupposition}.

\section{Parsons's refinement: four inferential modes}
\label{sec:p13-supposition-refinement}

Parsons proposes subdividing distribution by whether the entire
conjunction of instances permits return. Wide distribution allows
that return; narrow distribution does not. The latter is not historical
immobility, since it still allows descent. Nor is it a fourth mode
defined through conjunctive-term grammar. Paul supplies a precedent for
wide distribution; Parsons presents narrow distribution as his own
refinement \cite[pp.~266--268]{Parsons2007FordhamSupposition}.

For a finite nonempty list, write \(F\) for the original proposition,
\(F_1,\ldots,F_n\) for its instances, \(C=\bigwedge_iF_i\), and
\(V=\bigvee_iF_i\). Entailments in the following table are schematic
tests with reference and exhaustiveness information held fixed. A
crossed entailment means that the corresponding implication is not
generally valid; it does not assert failure in every individual model.
\[
\begin{array}{l|cc}
\text{mode}&\text{descent}&\text{return}\\\hline
\mathrm{Det}&F\Rightarrow V&V\Rightarrow F\\
\mathrm{MC}&F\not\Rightarrow V&V\Rightarrow F\\
\mathrm{W}&F\Rightarrow C&C\Rightarrow F\\
\mathrm{N}&F\Rightarrow C&C\not\Rightarrow F
\end{array}
\]
Here Det excludes the general conjunctive descent, W excludes ascent
from a single instance, and N excludes return even from the whole
conjunction. These qualifications separate the modes in the schematic
setting; singleton or logically fixed occurrences are not counterexamples
to the intended classification.

\begin{hbtheorem}{thm:p13-supposition-four-modes}
 {proved=annex:p13.supposition.four-modes}
 {The four patterns in the categorical core}
With nonempty subject and predicate ranges and independent schematic
term occurrences, the declared categorical clauses have the following
patterns for subject and predicate:
\[
 A:(\mathrm W,\mathrm{MC}),\qquad
 I:(\mathrm{Det},\mathrm{Det}),\qquad
 E:(\mathrm W,\mathrm W),\qquad
 O:(\mathrm{Det},\mathrm N).
\]
In the finite instance comparison, simultaneous negation of the
proposition and its instances interchanges Det with W and MC with N.
\end{hbtheorem}

The distinction explains the double-distribution problem. Negating a
wide universal produces a wide existential, hence determinacy. Negating
a narrow universal leaves an existential under a different universal,
hence mere confusion. The common name ``distributed'' previously hid
that difference. The more symmetric classification treats the modes
as wide or narrow, existential or universal, instead of treating every
non-determinate mode as one undifferentiated confusion
\cite[pp.~268--270]{Parsons2007FordhamSupposition}.

\subsection{The revised rules and their exact scope}
\label{sec:p13-supposition-prefix}

Normalize a categorical with fixed, independently ranging main terms
by replacing a negative quantifier with a universal followed by
negation, making indefinite signs explicit, and pushing each negation
rightwards across denoting phrases while dualizing \emph{every} and
\emph{some}. Singular phrases commute with propositional negation in
the nonempty comparison. Cancel adjacent double negations. The order
of the main phrases is preserved: this procedure never exchanges a
universal with an existential. The result is the designated prenex
convert, not a claim that all equivalent prenex formulas have one
unique ordering. The equivalences here require the nonempty ranges
already stipulated; empty-range affirmative truth must still be treated
by its separate clauses
\cite[pp.~270--273]{Parsons2007FordhamSupposition}.

A normalized main term has universal or existential import according
to its final sign. It is wide if no earlier main term has the opposite
sign, and narrow otherwise. This is a syntactic measure of the
quantificational force available in the whole context. It need not be
the widest scope obtainable from a special matrix by accidental logical
equivalences. Under the identification
\[
\mathrm{Det}=(\exists,\text{wide}),\quad
\mathrm{MC}=(\exists,\text{narrow}),\quad
\mathrm W=(\forall,\text{wide}),\quad
\mathrm N=(\forall,\text{narrow}),
\]
the four operations affect a following occurrence as follows. They
are applied from the innermost, rightmost sign outward.
\[
\begin{array}{c|cccc}
\text{new outer sign}&\mathrm{Det}&\mathrm{MC}&\mathrm W&\mathrm N\\\hline
\text{some}&\mathrm{Det}&\mathrm{MC}&\mathrm N&\mathrm N\\
\text{every}&\mathrm{MC}&\mathrm{MC}&\mathrm W&\mathrm N\\
\text{no}&\mathrm W&\mathrm N&\mathrm{MC}&\mathrm{MC}\\
\text{not}&\mathrm W&\mathrm N&\mathrm{Det}&\mathrm{MC}
\end{array}
\]
The new term itself is Det under \emph{some}, W under \emph{every}
or \emph{no}. These are Parsons's refined rules stated as a transition
table; they are not attributed to any medieval author.

\begin{hbtheorem}{thm:p13-supposition-prefix-rules}
 {proved=annex:p13.supposition.prefix-rules}
 {The revised causes compute the designated prefix classification}
For a categorical prefix made from fixed, nonempty independent ranges,
\emph{some}, \emph{every}, \emph{no}, and propositional \emph{not},
the transition table assigns to each occurrence exactly the sign and
wide/narrow status of its order-preserving prenex convert.
\end{hbtheorem}

The theorem establishes agreement of two explicitly declared
classifications. It does not settle arbitrary belief contexts, dependent
ranges, or modes of a term inside another term. In particular, a parasitic
main term still has a prefix position and receives a causal label,
although its unfixed range prevents the ordinary exhaustive-instance
test. That distinction is essential to the source's conclusion
\cite[pp.~272--273]{Parsons2007FordhamSupposition}.

\subsection{The nine possibilities and the limits of extension}
\label{sec:p13-supposition-other-modes}

With finite nonempty instances, descent can be conjunctive, merely
disjunctive, or absent. Ascent can come from a disjunction, merely from
a conjunction, or from neither. This gives nine formal combinations.
Using C and V from above, the source's chart is:
\[
\begin{array}{c|ccc}
 &\text{return from V}&\text{return from C only}&\text{no return}\\\hline
\text{descent to C}&[5]&\mathrm W&\mathrm N\\
\text{descent to V only}&\mathrm{Det}&[6]&[7]\\
\text{no descent}&\mathrm{MC}&[8]&[9]
\end{array}
\]
Possibility [5] would make every individual instance equivalent to the
original; its availability under the intended nondegenerate formal
tests is disputed. [6] is illustrated by an invented many-or-all
quantifier; [7] by some-but-not-all; [8] by every-or-no; and [9] by
a report that Maria said some donkey was spotted. The last can be true
without any donkey being spotted, while even all donkeys being spotted
does not make the report true. For [7] and [8], treating the denoting
phrase as negative changes which negation is retained in an instance.
These possibilities explain what the limited ordinary quantifier
vocabulary excludes; the source does not claim that medieval authors
developed this whole nine-case theory
\cite[pp.~273--274]{Parsons2007FordhamSupposition}.

Non-main terms raise a different extension problem. A determinate
relative-clause occurrence inside a determinate outer complex can have
determinate import globally: in \emph{Some man sees a horse which
some donkey sees}, the two witnesses can be chosen together. Replacing
the outer \emph{some man} by \emph{every man} can leave the donkey
witness dependent on the man, producing mere confusion. Intersective
adjectives and participles obtained from relative clauses provide further
instances of this analysis. This does not yield a composition table
for arbitrary pairs of local modes
\cite[pp.~274--276]{Parsons2007FordhamSupposition}.

\begin{hbtheorem}{thm:p13-supposition-relative-check}
 {proved=annex:p13.supposition.relative-check}
 {Relative-clause witnesses and a correction to the closing examples}
In finite relational comparisons with nonempty independent ranges,
the donkey occurrence in \emph{Some man sees a horse which some
donkey sees} has the determinate pattern; with \emph{every man}
it has the merely confused pattern. The occurrence in \emph{Some
horse is an animal that sees every donkey} instead has narrow
distribution, not mere confusion.
\end{hbtheorem}

The last conclusion corrects the active classification of the source's
second closing example, which prints ``merely confused'' and groups
it with examples whose outer complex has narrow distribution. Here
the outer predicate is existential, and the donkey term is a universal
within it. One horse seeing all donkeys entails a possibly different
horse for each donkey; the converse fails. The held manuscript's other
closing examples vary identity, seeing, and negation inside a complex
term, so their printed list cannot be used as a checked mode-composition
table. Parsons himself leaves the general algorithm for non-main terms
unconstructed. The proof just given repairs the explicit relational
case; it does not supply that missing general theory
\cite[p.~276]{Parsons2007FordhamSupposition}.

The resulting picture retains three levels. A term's signification
provides what it can stand for; the proposition's tense, modality,
and semantic context determine its actual supposition; ascent and
descent describe the inferences permitted at a particular occurrence.
The later theory makes that inferential role global. Parsons's modern
refinement explains a substantial extensional fragment, while dependent
ranges, opaque contexts, and complex-term composition keep the stated
boundaries visible. His manuscript's bibliography identifies the
underlying editions, translations, and interpretive disputes; their
presence here as mediated sources does not count as independent
collation of those works or of the distinct final 2008 chapter
\cite[pp.~276--280]{Parsons2007FordhamSupposition}.
